\chapter{Estado del arte}\label{cap:edoarte}

En este capítulo se presenta un panorama del estado del arte y el posicionamiento de esta tesis.

\section{Líneas base de recuperación: de BM25 a los modelos neuronales modernos}\label{sec:edoarte-lineas-base}

\subsection{Modelos léxicos: BM25 y variantes modernas}\label{subsec:edoarte-bm25}
% - Okapi BM25 como línea base omnipresente en benchmarks
% - Variantes (BM25+, BM25L): qué deficiencia corrige cada una
% - Infraestructura reproducible: Anserini/Pyserini (mención)

\subsection{Codificadores duales multilingües}\label{subsec:edoarte-duales}
% - Línea evolutiva: mDPR → mContriever → mE5 → mE5-instruct
% - Generación reciente: Qwen3-Embedding, jina-embeddings-v5
% - Qué reportan en multilingüe y en español (MIRACL, MTEB)
% - Cierre: criterio de selección de los modelos de esta tesis

\subsection{Modelos multifuncionales y dispersos aprendidos (BGE-M3, SPLADE)}\label{subsec:edoarte-multifuncionales}
% - SPLADE: v1 → v3, qué mejora cada versión y resultados reportados
% - BGE-M3: denso + disperso + multi-vector en un solo modelo
% - Anticipación de la fusión interna (enlaza con §3.3.3)

\subsection{Interacción tardía (de ColBERT a jina-colbert-v2)}\label{subsec:edoarte-tardia}
% - ColBERT: la propuesta original de interacción tardía
% - ColBERTv2: compresión y destilación
% - jina-colbert-v2: variante multilingüe; resultados reportados

\section{Reordenamiento neuronal en la literatura}\label{sec:edoarte-reranking}

\subsection{Cross-encoders (monoBERT, monoT5, bge-reranker-v2-m3)}\label{subsec:edoarte-crossencoders}
% - monoBERT y monoT5: reranking punto a punto con Transformers
% - bge-reranker-v2-m3: el cross-encoder multilingüe de esta tesis
% - Efectividad reportada y costo computacional

\subsection{Rerankers listwise con LLMs (RankGPT, jina-reranker-v3)}\label{subsec:edoarte-listwise}
% - RankGPT/RankZephyr: reranking como generación de ordenaciones
% - jina-reranker-v3: el listwise de esta tesis
% - Cierre de la sección: trade-offs efectividad/eficiencia en la cascada

\section{Fusión de rangos en la literatura}\label{sec:edoarte-fusion}

\subsection{Los métodos clásicos}\label{subsec:edoarte-clasicos}
% - Fox & Shaw 1994: CombSUM/CombMNZ en TREC
% - Aslam & Montague 2001: Borda-fuse (y Bayes-fuse)
% - Montague & Aslam 2002: Condorcet-fuse
% - Cormack et al. 2009: RRF
% - Bailey et al. 2017: RBC
% - Hallazgo transversal: la fusión suele superar a los componentes individuales

\subsection{Fusión léxico-semántica en la era neuronal}\label{subsec:edoarte-era-neuronal}
% - Karpukhin et al. 2020 (DPR): el denso supera a BM25, pero la combinación ayuda
% - Luan et al. 2021: combinaciones sparse/dense
% - Bruch et al. 2023 y Lin et al. 2021: análisis de cuándo aporta la fusión
% - Síntesis crítica: condiciones bajo las que la fusión léxico-semántica ayuda y cuándo no

\subsection{Fusión interna del modelo vs. fusión externa de listas (BGE-M3)}\label{subsec:edoarte-fusion-interna}
% - BGE-M3 como caso de estudio de fusión interna
% - Qué pierde la fusión interna frente a fusionar listas de sistemas especializados
% - Motivación directa de las preguntas de investigación de la tesis

\section{Benchmarks de IR y el español}\label{sec:edoarte-benchmarks}

\subsection{De TREC a MS MARCO y BEIR}\label{subsec:edoarte-trec-marco-beir}
% - TREC: pooling y juicios incompletos
% - MS MARCO: escala web
% - BEIR: evaluación zero-shot heterogénea

\subsection{Benchmarks multilingües y particularidades del español}\label{subsec:edoarte-espanol}
% - MIRACL, Mr. TyDi, MTEB Multilingual: cobertura del español
% - Morfología del español y sus efectos en modelos léxicos y tokenización
% - Backbones en español (BETO, RoBERTa-bne, MARIA): mención
% - Translationese: mMARCO-es vs. MIRACL-es vs. datos nativos

\subsection{MessIRve}\label{subsec:edoarte-messirve}
% - Descripción: origen, tamaño, variedades dialectales del español (incl. rioplatense)
% - Qué lo distingue: consultas nativas, no traducidas
% - Resultados previos reportados (incluidas las líneas base propietarias)

\section{Posicionamiento de esta tesis}\label{sec:edoarte-posicionamiento}
% - Gap 1: la fusión se evalúa sobre todo en inglés; falta evidencia sistemática en español
% - Gap 2: no hay comparación sistemática de algoritmos de fusión frente a rerankers modernos
% - Gap 3: abiertos vs. propietarios sin resolver en español
% - Puente al cap. 4: los cuatro experimentos responden a estos gaps
